\section{数据集与任务定义}

\subsection{XJTU-SY 数据与划分}

XJTU-SY 是轴承 run-to-failure 数据集。每个轴承在固定工况下从较健康状态运行到失效附近，原始数据以连续的振动快照组织；本项目的数据加载器把每个 \config{acc} 文件视为一个快照样本，并在索引中保留数据集、轴承、工况、样本序号与时间顺序。

本轮特征分析没有只依赖一个随机划分，而是使用三层检查：

\begin{itemize}[leftmargin=2em]
  \item 主划分 \config{xjtu_bearing_index_split}：同一工况下的前 3 个轴承进入训练集，第 4 个轴承进入验证集，第 5 个轴承进入测试集。这个划分用于回答“在覆盖三种工况的主训练设置下，哪些特征最有用”。
  \item 工况内 leave-one-bearing-out：分别在 35Hz12kN、37.5Hz11kN、40Hz10kN 内训练/验证/测试。这个划分用于回答“同一工况下，主结论是否稳定”。
  \item 跨工况划分 \config{xjtu_cross_condition}：训练集为 35Hz12kN，验证集为 37.5Hz11kN，测试集为 40Hz10kN。这个划分用于回答“特征排序是否只记住了某个工况的分布”。
\end{itemize}

三层检查是必要的。XJTU-SY 的工况差异会改变振动幅值、频谱位置和噪声水平；如果只看主划分，某些水平通道或频域特征可能看起来很强，但跨工况后会因为均值漂移或方差变化而降级。

\begin{table}[htbp]
\centering
\caption{XJTU-SY 分析划分}
\label{tab:xjtu_split_summary}
\scriptsize
\begin{tabularx}{\textwidth}{lYYYY}
\toprule
场景 & Train bearings & Val bearings & Test bearings & 用途 \\
\midrule
main & B1\_1--B1\_3; B2\_1--B2\_3; B3\_1--B3\_3 & B1\_4; B2\_4; B3\_4 & B1\_5; B2\_5; B3\_5 & 主三任务 ranking，7032/1679/505 samples。\\
C1 LOO & B1\_1--B1\_3 & B1\_4 & B1\_5 & 35Hz12kN 工况内稳定性，442/122/52 samples。\\
C2 LOO & B2\_1--B2\_3 & B2\_4 & B2\_5 & 37.5Hz11kN 工况内稳定性，1185/42/339 samples。\\
C3 LOO & B3\_1--B3\_3 & B3\_4 & B3\_5 & 40Hz10kN 工况内稳定性，5405/1515/114 samples。\\
cross & B1\_1--B1\_5 & B2\_1--B2\_5 & B3\_1--B3\_5 & train=35Hz12kN, val=37.5Hz11kN, test=40Hz10kN，616/1566/7034 samples。\\
\bottomrule
\end{tabularx}
\vspace{0.3em}
\begin{flushleft}
\scriptsize 注：B1\_k 表示 Bearing1\_k，B2\_k 表示 Bearing2\_k，B3\_k 表示 Bearing3\_k。
\end{flushleft}
\end{table}


\subsection{PHM2012 数据与划分}

PHM2012 也是轴承 run-to-failure 数据集。当前框架按官方数据目录构建索引，并使用显式的 \config{phm2012_official} 划分。训练集、验证集和测试集按轴承隔离，保证同一轴承的时间序列不会同时出现在不同 split 中。

PHM2012 的本轮分析重点不是故障类型分类，而是退化建模。其原因有两点。第一，本阶段要统一两个数据集的三任务主线：RUL、Health State、Early Fault Detection。第二，Health State 和 Early Fault 不是人工逐点标注的真实类别，而是由健康指标 HI 和首次预测时间 FPT 规则构造的 pseudo-label；它们更适合用于退化阶段解释和早期预警分析，而不是直接等价为官方故障类型标签。

\begin{table}[htbp]
\centering
\caption{PHM2012 official split}
\label{tab:phm2012_split_summary}
\scriptsize
\begin{tabularx}{\textwidth}{lYr}
\toprule
Split & Bearings & Samples \\
\midrule
Train & Bearing1\_1, Bearing1\_2, Bearing2\_1, Bearing2\_2, Bearing3\_1, Bearing3\_2 & 7534 \\
Val & Bearing1\_3, Bearing2\_3 & 4330 \\
Test & Bearing1\_4, Bearing1\_5, Bearing1\_6, Bearing1\_7, Bearing2\_4, Bearing2\_5, Bearing2\_6, Bearing2\_7, Bearing3\_3 & 13025 \\
\bottomrule
\end{tabularx}
\end{table}


\subsection{三任务定义}

本报告统一分析三个任务。

\begin{description}[leftmargin=2.4em,style=nextline]
  \item[RUL]
  RUL 使用 \config{piecewise_rul_norm} 作为目标列。它表达样本距离失效终点的归一化寿命位置，适合用于回归、排序或趋势建模。特征分析时主要看特征是否随寿命推进呈稳定趋势。

  \item[Health State]
  Health State 使用 \config{health_state_id}。它是由 HI/FPT 逻辑派生的 4 类 pseudo-label，表示 healthy、slight、moderate、severe 等退化阶段。它适合看不同退化阶段之间的特征分布是否可分。

  \item[Early Fault Detection]
  Early Fault 使用 \config{early_fault}。它是 FPT 前后构造的二分类 pseudo-label，用于看某个特征在首次可预测退化点附近是否出现明显变化。
\end{description}

\subsection{为什么本轮不做 Fault Type}

Fault Type 任务需要稳定且可比的故障类型标签，例如 outer race、inner race、ball fault 等明确类别。本轮两个数据集的统一目标是 run-to-failure 退化过程分析，而不是故障类别识别。为了避免把不同来源、不同粒度的标注混在一起，当前三任务主线明确排除 Fault Type。

这也影响推荐特征的解释边界：本报告推荐的特征适合 RUL、Health State 和 Early Fault Detection；它们不自动等价为 fault type classifier 的最优特征。后续若进入故障类型识别，需要重新定义标签、划分和分析指标。

\begin{table}[htbp]
\centering
\caption{特征分析步骤与结论}
\label{tab:analysis_runs}
\small
\begin{tabularx}{\textwidth}{llY}
\toprule
步骤 & 对象 & 结论 \\
\midrule
Step F & XJTU-SY \config{manual_basic} 主划分 & RUL 和健康状态主要由幅值/能量类时域特征支撑。\\
Step G & XJTU-SY \config{manual_tsfresh_basic} & full-size tsfresh 被阻塞并延后，当前主线继续 \config{manual_basic}。\\
Step H & XJTU-SY 工况内分析 & RUL 稳定，Early Fault 最工况敏感。\\
Step I & XJTU-SY 跨工况分析 & RUL 幅值特征保留，HealthState 更偏 magnitude 稳定，Early Fault 需降级为条件性结论。\\
Step J & PHM2012 \config{manual_basic} & 三任务均由 horizontal/magnitude 幅值特征主导。\\
Step K & PHM2012 \config{manual_tsfresh_basic} & 能跑通，但 top tsfresh 特征主要重复 RMS/std/variance 等手工统计量。\\
\bottomrule
\end{tabularx}
\end{table}

